购买系统的在线链路按分区保存受理、付款、库存决定和事件责任。
批量对账从这些保存域取得确定范围的材料，把分散记录按业务关系汇合起来。
它可以发现原付款记录 $C$ 已有平台收款而权威结果尚未补齐，
也可以检查已付款购买是否形成原订单与事件 $O/EV$、
消费者是否保存原接收记录与履约效果 $I/F$。
作业执行者能够更换，参加这次计算的数据范围与记录关联则由作业清单保存。

这类工作把两种生命周期放在同一幅系统图里。
业务权威保存原 $S$ 到 $P/C/T$ 的绑定，以及跨分区决定、未决责任和原结果。
计算运行时负责把程序放到可用机器上，取得输入，组织依赖，回收执行资源。
批处理适合扫描大量固定材料，Ray 适合在计算过程中继续展开任务，
Actor 可以长期持有某分区的执行缓存，事件函数适合按事件取得一段执行时间。
它们通过同一恢复入口读原责任、查询原平台操作并回写结果，
计算进程的寿命因而可以短于一笔付款的恢复时间。
图~\ref{fig:ch20-runtime-place} 先给出它们在完整购买系统中的位置。

\CTwentyFigure{ch20-22}{执行运行时承载计算，持久保存域连接整条业务责任链}{fig:ch20-runtime-place}

对账输入包含权威图像、平台账本和消费账本时，三个保存域通常没有共同的全局快照。
作业清单分别保存来源、版本或截止位置，以及取得材料的条件。
假定权威图像停在应用索引 36，平台导出截止于稍后时刻，
平台多出的一笔收款可能正是权威尚未观察到的结果。
对账把它识别为需要核实的差异，再按原 $C/T$ 查询当前权威。
恢复者若取得平台原结果 $X_1$，便将这份结果补回原记录。
不同导出位置造成的观察缺口，经过查询才能转成确定的处置依据。

\CTwentyFigure{ch20-23}{来源截止位置不同，差异先表示观察缺口，再由原操作查询确定}{fig:ch20-input-cuts}

MapReduce 或 Ray 读取固定输入、计算差异并发布结果，
支付修复工作者根据差异回到原责任。
函数执行环境可以承载一次查询、回填或消息处理，
长期未决的 $C/T$、跨分区事务 $TX$ 与退款占额保存在各自权威域。
图~\ref{fig:ch20-batch-recovery} 中，差异清单连接计算输出与业务恢复入口。

\CTwentyFigure{ch20-01}{批量计算从固定材料形成差异，再回到原责任恢复}{fig:ch20-batch-recovery}

\subsection{作业身份先固定输入范围}

一次对账作业 $J_1$ 保存完整输入清单：分区 $b_1$ 的图像版本、
平台账本导出版本、消费账本截止位置、解析格式与计算代码版本。
作业重试读取同一份清单。
运营人员改用下一天账本或改变差异规则时，创建新的作业版本，
并保存它与前一结果的关系。
若名称只写昨日对账，时区调整、迟到数据和规则变化就会挤进同一个名称，
恢复者无法从名称判断两次执行应当重用哪些材料。

作业内部再区分逻辑任务与执行尝试。
逻辑任务 $M_3$ 表示处理 $J_1$ 中指定的一段输入，
尝试 \code{M3-a1} 与 \code{M3-a2} 表示不同进程对这段输入的执行。
任务重调度改变尝试身份，输入区间与输出职责保持相同。
最终输出身份由 $J_1$ 及输出分区确定，
尝试的临时文件名只定位尚未发布的候选字节。

这些身份服务于不同关联。
$J_1$ 关联输入与整个结果集合，$M_3$ 关联一份逻辑计算，
\code{a1} 关联执行位置和故障诊断。
购买身份 $S$、原 $C$ 和平台键则跨越多个作业存在。
$J_1$ 与 $J_2$ 都发现同一个未决 $C$ 时，
两个作业可以各自产生一条差异记录，后续修复通过原 $C$ 汇合到同一次付款责任。

\CTwentyFigure{ch20-02}{作业、逻辑任务、执行尝试与业务身份分别定位}{fig:ch20-identities}

对账记录保留完整关联，才能判断观察到的是哪笔款。
权威记录携带完整 $S$、原输入、$P/C/T$、原 \code{scope} 与 \code{pay-C}，
以及已经保存的受理事实和原结果。
平台记录保留原平台身份、金额、币种和原收款结果。
汇合后逐项比较这些关联，才能区分同一笔款的重复观察、
输入冲突以及真正不同的购买。
只比较成功数为一，会把这些不同历史压成同一个统计值。

假定本次固定材料中，Alice 的账户 1 有一笔 3000 分购买，
原 $C_1$ 关联平台结果 $X_1$，另有一笔因库存不足终结的购买。
map 阶段把权威记录和平台记录转换成带来源的键值记录，
以原平台 \code{scope} 与 \code{pay-C} 组成中间键。
reduce 阶段在同一个键下取得两类记录，
核对完整输入和原结果，再输出匹配、待查询或冲突。
本次计算使用固定材料里的金额，当前商品报价没有参加这笔旧购买的输入绑定。

\subsection{MapReduce 怎样分配计算}

MapReduce 将大规模数据处理拆成两类操作。
map 逐条读取输入，提取后续需要汇合的键和内容。
reduce 取得同一个键下的内容集合，执行分组核对或汇总。
运行时在两者之间完成分区、传输、排序和归并，
应用因此可以围绕每条记录及同键记录编写计算，
把大量执行位置的管理交给运行时~\cite{mapreduce-design}。

输入先切成多个 split，即可独立读取的逻辑输入片段。
它包含文件范围和读取规则，切分还要尊重记录边界。
输入片段与底层存储块各有用途：前者分配计算，
后者保存字节和副本，一个任务可以根据读取规则涉及一个或多个块。
map 产生的中间键经过分区函数分配给 reduce 任务，
相同键的记录进入同一 reduce 分区。
多个不同键也可能进入同一分区，reduce 因而还要按键排序和分组。
图~\ref{fig:ch20-key-flow} 展示两份来源怎样沿平台身份汇合。

\CTwentyFigure{ch20-24}{带来源的记录经过分区与按键归并，形成同一付款的比较集合}{fig:ch20-key-flow}

在 Google MapReduce 的原设计中，master 是作业的中央调度进程，
保存任务状态、执行位置和中间文件位置。
worker 是接受任务的执行进程，负责读取、计算和实际数据传输。
map 先把中间记录放入内存缓冲，再分批写到本地磁盘，
每个 reduce 分区对应一个逻辑片段。
worker 向 master 报告这些片段的位置和大小，
master 再把位置传给相关 reduce worker。
reduce 直接从 map worker 取得片段，进行排序、归并和分组。
中央调度者传播位置，大块中间字节在 worker 之间流动。

这与第十九章中元数据服务定位块、副本直接传输字节的分工相近，
但 map 中间文件的保存责任较短。
固定输入已保存在可靠存储中，本地片段丢失后可以重算。
最终 reduce 输出保存到可靠文件系统后，
后续作业便可把它作为新的固定输入。
输入、本地中间数据与已发布结果由此有不同的恢复入口。

\CTwentyFigure{ch20-03}{master 传播位置，map 与 reduce 之间直接传输中间数据}{fig:ch20-map-runtime}

master 为每个任务保存空闲、执行中或已完成状态，
并保存执行中及已完成任务所在的 worker。
map 完成后还要保存各 reduce 片段的位置。
某个 worker 失去响应时，master 将其正在执行的任务重新置为空闲，
也重新调度那些只在该 worker 本地保留输出的已完成 map。
已经写入全局文件系统的 reduce 输出，
按该存储系统的保存条件继续存在~\cite{mapreduce-design}。

设 $M_3$ 在 worker A 生成片段，$R_1$ 尚未取得它，A 随后退出。
master 把 $M_3$ 分给 B，B 读取同一输入区间，
重新生成片段，$R_1$ 改从 B 读取。
$R_2$ 若已在故障前读走 A 的片段，可以继续计算。
对确定性 map，两份片段表达相同记录。
若 map 读取实时价格或其他变化数据，$R_1$ 与 $R_2$ 可能结合不同观察，
各自完成的输出分区便未必属于同一次共同观察。
作业清单把所需版本固定下来后，重算才有可重复取得的输入。

\CTwentyFigure{ch20-25}{同一 map 的位置改变，已读取和未读取的 reduce 分别继续}{fig:ch20-map-recovery}

master 自身退出则涉及它保存的调度状态。
原设计说明可以定期保存这些状态的检查点，
但所描述的实现选择在 master 失败时终止作业，
由调用方决定是否重新执行\footnote{原文为
``our current implementation aborts the MapReduce computation if the master fails''。
这一实现选择见集中书目所列 MapReduce 设计说明的故障处理部分。}。
调用方重新启动时，需要按照作业身份辨认输入和已有输出，
再建立新的调度状态，才能决定哪些成果可以继续使用。

\subsection{数据分布决定计算与恢复成本}

分区函数同时影响数据汇合和负载。
按完整平台身份分区，可以让同一笔付款的所有记录进入同一个 reduce。
若只按账户分区，一个大商户可能集中大量记录，形成长尾任务。
可以先按细粒度业务键汇合，再做账户级汇总，
第二层汇总保存第一层已经核对的付款关联，
便能在总额出现差异时回到具体的原 $C$。

combiner 是 map 侧可选的局部合并步骤，
用于减少送往 reduce 的中间记录。
可分段合并的计算需要保持相应的结合、交换等运算性质，
使运行时改变局部合并次数和分组方式后，最终含义一致。
精确整数分单位的金额总和可以逐段相加，
但同一 $C$ 的重复平台记录要先按身份和原内容核对。
若先把每次观察的金额都相加，重投的同一条记录就会进入两次总额。

调度还要考虑输入位置。
把 map 放在输入副本附近，可以减少跨机传输。
把任务切得比 worker 数更多，可以将一个故障 worker 的任务分散给多个空闲 worker。
过细切分增加调度、文件和元数据开销，
过粗切分扩大单次重算与长尾影响。
任务粒度决定节点退出后要重新读取多少输入、
重新计算多少记录，以及重新传输多少中间字节。
